\documentclass[10pt]{letter}
\usepackage[a4paper, margin=0.85in]{geometry}
\usepackage[T1]{fontenc}
\setlength{\parskip}{5pt}
\signature{Harsh Bordekar \\ harshbordekar@gmail.com}
\begin{document}
\begin{letter}{Airbus Aerostructures GmbH \\ Augsburg}
\opening{Dear Hiring Team,}

I am writing to apply for your development-engineer position in composite manufacturing technology with a focus on Automated Fibre Placement (AFP) at your Augsburg site. I am a composite research and simulation engineer with over seven years in CFRP and high-fidelity FEM, and, importantly for this role, my PhD is built directly on AFP-processed material.

At DLR and TU Braunschweig I am developing a data-driven multiscale simulation framework for AFP-processed CFRP. Using microscopy-derived microstructure data, it captures AFP-induced manufacturing nonlinearities and defects and propagates their effect from ply and laminate up to structural scale, extending to leakage prediction in AFP-manufactured hydrogen pressure tanks. This has given me a strong working understanding of how AFP process features and manufacturing variability translate into simulation and into predicted part quality, which maps closely to your tasks in process modelling, process optimisation, and quality prediction. I work daily in ABAQUS, including thermal and structural simulation on HPC.

Beyond process modelling, I evaluate composite material systems and effective properties through computational homogenisation, enabling fast material screening, and I quantify manufacturing-induced defects such as fibre waviness and porosity from microscopy and CT data. This directly supports the assessment of new material systems and manufacturing technologies for quality and robust series production that your team drives.

I am also used to the research-to-production setting this role sits in. I contribute to and coordinate work packages within collaborative, publicly-funded research projects, interface with partner institutes and universities, support knowledge transfer from research into application, and document results to aerospace reporting standards and scientific practice.

To be transparent about fit: my AFP experience is on the simulation and material side rather than machine operation, and I am keen to build hands-on AFP-machine and CATIA composite-design experience and to extend into thermoplastic joining. My German is currently B1 and I am actively working toward B2. I am strongly motivated by Airbus Aerostructures' role in preparing the composite airframe of the next aircraft generation, and I would welcome the chance to bring my simulation and materials background into your technology-development team.

Thank you for your consideration. I look forward to discussing how I can contribute.

\closing{Sincerely,}
\end{letter}
\end{document}